\section{Por qué la terminal vuelve a ser la entrada de la IA}

\subsection{La terminal no es una interfaz antigua, sino la capa de ejecución mínima y componible}

Si consideramos el IDE como el centro de edición de código y el navegador como el centro de consumo de información, entonces la terminal se parece más a un «plano de ejecución» (execution plane) del mundo de la ingeniería. Tiene tres ventajas difíciles de sustituir:

\begin{itemize}
    \item \textbf{Orquestable de manera natural}: los comandos, las tuberías, las variables de entorno y los scripts son en sí mismos bloques componibles
    \item \textbf{Amplia cobertura}: la compilación, las pruebas, el despliegue, la resolución de fallas y el procesamiento de datos pueden llevarse a la línea de comandos
    \item \textbf{Facilidad para volverse agencial}: un Agent puede convertir sus acciones en secuencias de comandos reproducibles, auditables y reintentables
\end{itemize}

\begin{knowledgebox}{El punto de encaje entre la terminal y los Agent}
Para un Agent, la terminal no es solo una interfaz de caja negra, sino una interfaz del sistema operativo con entradas, salidas, límites de permisos y señales de error claros. En comparación con una interacción completamente gráfica, la línea de comandos se presta mejor a que el modelo realice el ciclo de planear, ejecutar, observar y corregir.
\end{knowledgebox}

\subsection{El salto de interacción de «autocomplete» a «delegate»}

Los primeros productos de programación con IA se centraban sobre todo en el autocompletado (autocomplete) y la asistencia por chat (chat assist), pero Zach Lloyd enfatiza que el valor de la siguiente etapa proviene de la delegación (delegate):

\begin{enumerate}
    \item El desarrollador describe primero el objetivo, en lugar de describir la implementación línea por línea
    \item La herramienta genera un plan y muestra las acciones que podría tomar
    \item El usuario confirma en los puntos clave, en lugar de vigilar cada paso
    \item El sistema devuelve al desarrollador los resultados, los errores, los registros y el diff
\end{enumerate}

\begin{importantbox}{La dimensión de la competencia entre herramientas de desarrollo se desplaza hacia arriba}
La competencia ya no consiste solo en «quién autocompleta con mayor precisión», sino en «quién puede convertir de manera segura una intención de ingeniería de más alto nivel en acciones de ejecución concretas». La terminal es importante precisamente porque es lo más cercano a la ejecución real, y a la vez es más fácil de estandarizar que una interfaz gráfica.
\end{importantbox}

\subsection{Resumen de la sección}

La terminal vuelve a ser fundamental no porque sea algo retro, sino porque combina expresividad, capacidad de orquestación y capacidad de auditoría. Si las herramientas de IA quieren evolucionar de sistemas de sugerencia a sistemas de ejecución, la terminal será uno de los puntos de aterrizaje más naturales.
